% mixed-example.tex — template mixing prose, mathematics, and Lean 4 code.
% XeLaTeX (Menlo renders Lean's Unicode).  Compile: scripts/tex2pdf.sh scripts/mixed-example.tex
\documentclass[11pt]{article}
\usepackage[margin=1in]{geometry}
\usepackage{fontspec}
\setmonofont{Menlo}[Scale=0.85]
\usepackage{amsmath, amssymb}
\usepackage{fvextra}
\usepackage{xcolor}

\title{Prose, mathematics, and Lean 4 code together}
\author{ScottLean4 \textperiodcentered\ for Dana Scott}
\date{}

\begin{document}
\maketitle

This template shows the three kinds of content you will interleave: ordinary
\textbf{prose}, typeset \textbf{mathematics}, and blocks of \textbf{Lean~4 code}.

\medskip
Prose flows as usual, with inline math such as
$\forall n \in \mathbb{N},\; n + 0 = n$, and displayed math:
\[
  \sum_{k=0}^{n} k \;=\; \frac{n(n+1)}{2}, \qquad
  (a \wedge b) \to (b \wedge a).
\]

A block of Lean~4 code is set verbatim; its Unicode ($\mathbb{N}$, $\to$, $\forall$)
renders directly because we compile with XeLaTeX and a Unicode monospace font:

\begin{Verbatim}[fontsize=\small, frame=leftline, framerule=0.4pt, numbers=left]
theorem add_zero (n : ℕ) : n + 0 = n := by
  simp

theorem and_comm (a b : Prop) (h : a ∧ b) : b ∧ a :=
  ⟨h.right, h.left⟩
\end{Verbatim}

\noindent
And then prose resumes, and can refer back to the code above and to the display
equation. The pattern — \emph{paragraph, equation, code block, paragraph} —
repeats for as long as the document needs.

\end{document}
